\chapter{Einleitung}
In vielen Unternehmen entstehen Abläufe zunächst aus praktischer Notwendigkeit heraus, ohne dass sie von Anfang an systematisch dokumentiert oder modelliert werden \cite{Pospiech2014}. Solches informelles Vorgehen stößt jedoch an seine Grenzen, wenn die Komplexität zunimmt, die Anzahl der beteiligten Fachbereiche wächst und die Anforderungen an die Nachvollziehbarkeit steigen. Genau in dieser Lage befindet sich der Prozess Manufacture's Declared Value \acs{MDV} Verbrauch und Reichweite bei Mercedes-AMG. Obwohl an diesem Prozess verschiedene Fachbereiche wie Entwicklung, Zertifizierung und Vertrieb beteiligt sind und regulatorische Anforderungen eine hohe Verlässlichkeit verlangen, liegt bislang kein dokumentiertes oder modelliertes Prozessmodell vor \citeintern{amg2025mdvvorgehensweise}. Hierfür stellt das Geschäftsprozessmanagement (GPM) etablierte Methoden der Prozessaufnahme und -modellierung bereit, insbesondere die grafische Darstellung ``Business Process Model and Notation'' (BPMN) 2.0. Mit dieser lassen sich Aktivitäten, Verantwortlichkeiten und Datenflüsse einheitlich abbilden. Bei Mercedes-AMG werden unternehmensweite Prozesse auf der Plattform Signavio modelliert. Daraus ergibt sich damit folgende Forschungsfrage: Wie lässt sich der bisher undokumentierte MDV-Prozess für Verbrauch und Reichweite bei Mercedes-AMG mittels BPMN 2.0 in Signavio methodisch erfassen und modellieren, um Prozesstransparenz zu schaffen und den Ablauf in die übergeordnete Prozessarchitektur einzubinden? Zur Beantwortung dieser Frage werden im Folgenden zunächst die Problemstellung und die verwendete Methodik erläutert. Anschließend wird die Methodik im Rahmen des Transfers konkret auf den MDV-Prozess angewendet.
